% s05_specialisation (abridged version). Specialisation to the fibrewise torus.
\section{Specialisation to the fibrewise torus}
\label{sec:spe}

In this section the graph sum of \cref{thm:gen-main} is carried from generic weights to the fibrewise torus, by the ring map $\alpha_i\mapsto E/3$ on a subring of power series in $y$, with a denominator outside the prime $\primeq$ of $\TF$. No limit in the weights is taken, and the theorem of Teleman is not applied on $\TF$, where the cohomological field theory is not semisimple at $\cQ=0$.

\subsection{Statement}
\label{sec:spe-statement}

Let $\locO\coloneqq\ringA_{\primeq}$, with residue field $\Ky\coloneqq\bbK(y)$, and put
\[
\ringC\coloneqq\locO[Y]/(\Lhat(Y)),\qquad \redn\colon\locOal\fps{y}\longrightarrow\bbK\fps{y},\quad \alpha_i\mapsto E/3\ \text{ coefficientwise}.
\]
With $\Phi(\alpha,Y;z)$ the algebraic row \eqref{eq:row-Phi}, put $\Pn_0\coloneqq\Theta(z/\epsilon_1)\,\Theta(z/\epsilon_2)\,\Phi$, $\opD\coloneqq z\theta+Y-\frac{z}{2}\theta\log W$ and $\Qa\coloneqq\opD\Pn_0$. On $\TF$, at colour $i$, let $\Pb_a^{(i)}\coloneqq\Delta_i^{1/2}\RLR_a^{(i)}$ be the large-radius rows of \cref{def:set-rows} at $\cone=A$, multiplied by $\Delta_i^{1/2}$; thus $\Pb_0^{(i)}=\hb(z)\,\Srow(v_i,z)$. Let $\speEtab^{ab}\coloneqq\eta^{ab}|_{\alpha_i=E/3}$, and let the factors be those of \cref{thm:gen-main} for these rows:
\[
\Dilb_i(d)\coloneqq(-1)^d[z^{d-1}]\Pb_0^{(i)}\ \ (d\geq2),\qquad
\Edgeb_{ij}(\psi,w)\coloneqq\frac{\delta_{ij}\Delta_i-\sum_{a,b}\Pb_a^{(i)}(-\psi)\,\speEtab^{ab}\,\Pb_b^{(j)}(-w)}{\psi+w}.
\]

\begin{thm}[Specialisation; {\Lg{Theorem~5.1}}, condensed in parts \textup{(a)} and \textup{(b)}]\genusstatus{AI-proposed, unchecked}
\label{thm:spe-main}
Parts \textup{(c)}, \textup{(d)}, \textup{(e)} are for $g\geq 2$. The proof uses \cref{thm:gen-main}, the results of \cref{sec:row} and \cite{GP99}; part \textup{(b)} does not use the mirror theorem.
\begin{description}
\item[(a) Regularity of the factors, to all orders in $z$] $\ringC$ is free of rank $3$ over $\locO$ and local, with maximal ideal $\primeq\ringC$ and residue field $\bbL\coloneqq\bbK(t)$, $t^3=y$, via $\redC\colon Y\mapsto-(E/3)/(1+3t)$. The elements $Y$, $W$, $Y+\alpha_j$ and $\Lhat'(Y)=-W/Y$ are units of $\ringC$, and $\theta$ extends uniquely to $\ringC$. The series $\Pn_0$, $\Qa$, $\opD\Qa$ have $z$-coefficients in $\ringC$: no square root and no discriminant occurs in the denominators. For every colour $j$ at generic $\alpha$, $\Delta_j^{1/2}\Rm_{0j}=\iota_j(\Pn_0)$, where $\iota_j\colon\ringC\to\Kal\lser{y}$ is $Y\mapsto Y_j(y)$. In the pairing of the rows the function $A$ cancels, and $\Edge_{ij}(a,b)$ lies in $(\iota_i\otimes\iota_j)(\ringC\otimes_{\locO}\ringC)+F(\iota_i\otimes\iota_j)(\ringC\otimes_{\locO}\ringC)$, linear in $F$, with $\de_{\sfF}\Edge_{ij}(\psi,w)=3\epsilon_1\epsilon_2\,\xi\,\Qa^{(i)}(-\psi)\,\Qa^{(j)}(-w)$. The symplectic identity and the delta form of the edge hold with $\sfF$ free. The map $\redC$ carries all of this to the data on $\TF$: at each $z$-order, $\redC(\Phi)=\Theta(-z/E)\,\Srow(v,z)$ and $\redC(\Pn_0)=\hb\,\Srow$.
\item[(b) The Gromov--Witten side] The statements of \cref{prop:gen-poly} hold. The constant term is as follows: if $H'$ is the non-equivariant hyperplane class of $\bbP^2$, then for $g\geq2$
\[
\GW_{g,0}^{\TF}=\int_{\Mbar_g}\big[9\lambda_g\lambda_{g-2}\mathrm{P}(0)-3\lambda_g\lambda_{g-1}\mathrm{P}'(0)\big],\qquad
\mathrm{P}(H')\coloneqq\frac{\Lamv(E-3H')\Lamv(-\epsilon_1)\Lamv(-\epsilon_2)}{(E-3H')\,\epsilon_1\epsilon_2}.
\]
For $g=2$ this is $E^2/640$, and for $g=3$ it is $E^4/161280$. No closed form is claimed for $g\geq 4$.
\item[(c) Generic $\alpha$, algebraic form] For each stable graph $\Gamma$ of genus $g$ without legs and each decoration, the term of the graph sum \eqref{eq:gen-graphsum} at a colouring $\col$ is $(\prod_v\iota_{\col_v})$ applied to an element $\Tsum_{\Gamma,\mathrm{dec}}$ of $\ringC^{\otimes_{\locO}|\Ver(\Gamma)|}[\sfF]$. So
\[
\Tsum_g\coloneqq\sum_\Gamma|\Aut\Gamma|^{-1}\sum_{\mathrm{dec}}\big(\Tr_{\ringC/\locO}\big)^{\otimes|\Ver(\Gamma)|}\big(\Tsum_{\Gamma,\mathrm{dec}}\big)
\]
lies in $\locO[\sfF]$, $\deg_{\sfF}\Tsum_g\leq 3g-3$, and the identity \eqref{eq:gen-graphsum} reads $\sF_g^{\Tmax}=\Tsum_g(F)$ in $\Kal\fps{y}$, where the sum over colourings is the multi-trace.
\item[(d) The specialisation identity] In $\bbK\lser{y}$,
\begin{equation}
\label{eq:spe-identity}
\sF_g^{\TF}(y)=\Tsumb_g(F),
\end{equation}
where $\Tsumb_g(F)$ is the point-frame graph sum for the rows $\Pb$: the right side of \eqref{eq:gen-graphsum} with $\Dilb_i$, $\Edgeb_{ij}$ and the $\Delta_i$ of \eqref{eq:set-DeltaTF} in place of $\Dil_i$, $\Edge_{ij}$ and the $\Delta_i$ at generic weights. Each coloured term lies in $\bbK(\omega)[t^{\pm1},(1+3\omega^it)^{-1}:i][F]$ and has poles at $t=0$ in general. The colour sum lies in $\bbK[y^{\pm1},\xi^{-1}][F]$. The left side is a formal series in $\bbK\fps{y}$; its convergence near $y=0$ is a consequence of \eqref{eq:spe-identity}. Continuation to the universal cover of $\bbC\setminus\{0,-1/27\}$ means continuation of the right side, and it is not a statement about Gromov--Witten series at large $y$.
\item[(e) Consequences] $\sF_g^{\TF}$ lies in $\bbK[y^{\pm1},\xi^{-1}][F]$ with $\deg_F\leq 3g-3$. No finite generation statement is used, and neither is $A$, which cancels in graphs without legs. $F$ is transcendental over $\bbC(y)$, by an elementary monodromy argument, so the polynomial in $\sfF$ is unique, and $\Cf(\sF_g^{\TF})\coloneqq\Tsumb_g(\Fcf)$ is well defined (see \cref{def:fra-cf}). Colouring by colouring, the same holds over $\bbC(t)$.
\end{description}
\end{thm}

\subsection{The key points of the proof}

\emph{The ring $\ringC$ and the trace} \Lg{Lemma~5.3}. The polynomial $\Lhat$ has leading coefficient $\xi$, a unit of $\locO$, so $\ringC$ is free on $1,Y,Y^2$. Modulo $\primeq$ it becomes $(Y+E/3)^3+27yY^3$, and with $t\coloneqq-(Y+E/3)/(3Y)$ one gets $t^3=y$; so the quotient is the field $\bbL$, and $\ringC$ is local. At generic $\alpha$ the maps $\iota_0,\iota_1,\iota_2$ form an isomorphism after base change to $\Kal\lser{y}$, so $\Tr_{\ringC/\locO}=\sum_j\iota_j$, and for $T\in\ringC^{\otimes n}$ the multi-trace $\Tr^{\otimes n}(T)=\sum_{\col}(\iota_{\col_1}\otimes\dots\otimes\iota_{\col_n})(T)$ lies in $\locO$. The trace of a free module commutes with base change: the reduction of $\Tr^{\otimes n}(T)$ modulo $\primeq$ is $\Tr_{\bbL/\Ky}^{\otimes n}(\redC^{\otimes n}T)$, which after adjoining $\omega$ is the sum over colourings of $\redC^{\otimes n}T$ evaluated at $t_v=\omega^{\col_v}t$. So Galois invariance is replaced by the trace.

\emph{The rows and the pairing} \Lg{Lemmas~5.4, 5.5}. By \cref{prop:gen-osc}, with $\Delta_i=3\epsilon_1\epsilon_2W_i$ and $\Theta(x)\Theta(-x)=1$, one gets $\Delta_i^{1/2}\Rm_{0i}=\iota_i(\Pn_0)$, all branch choices cancelling. Note that $h_{\wts_i}$ and $\Rt$ are each singular on the diagonal $\alpha_0=\alpha_1=\alpha_2$, while their product is not: the Gamma part of $h_{\wts_i}$ is absorbed in $\Phi$. By the row recursion \eqref{eq:gen-rowrec}, $\Pn_1=\Qa/A$ and $\Pn_2=D^{-1}A^{-1}[\opD\Qa-zF\Qa-e_1(\alpha)\Ey\,\Qa]$. Inserting the relations of \cref{prop:row-modrel}, the terms in $A^{-2}$ cancel in the pairing, because $\eta^{11}=e_1(\alpha)\eta^{12}$ and $\eta^{22}=0$. The pairing $\langle\Pn(z),\Pn(w)\rangle$ becomes an explicit expression in $\Pn_0$, $\Qa$, $\opD\Qa$, linear in $\sfF$, whose $\sfF$-term carries the factor $(z+w)$ \Lg{eq.~(5.1)}. The edge is then a tensor in $\ringC\otimes_{\locO}\ringC[\sfF]$, and $|\Edg(\Gamma)|\leq3g-3$ bounds the degree in $\sfF$. This gives \textup{(a)} and \textup{(c)} at generic $\alpha$.

\emph{Reduction of the rows} \Lg{Lemma~5.6}. The identity $\redC(\Phi)=\Theta(-z/E)\,\Srow(v,z)$, with $v=3/(E+3Y)$, is where the closed form of \cref{sec:set} is identified with the canonical row. On $\TF$, in the coordinates $s_n=e_n(q_0^{1/3},q_1^{1/3},q_2^{1/3})$, the phase is $s_1^3-3s_1s_2+\kappa s_3-E\log s_3$ with $\kappa=Ev$, and it splits exactly. By coordinate invariance and factorisation (\cref{sec:row-sp}), the $(s_1,s_2)$-block contributes numbers $r^{\mathrm{W}}_l$, and the $s_3$-integral is the formal Stirling formula \eqref{eq:row-stirling} at $a=-E$, $\beta=-l-1$, which gives $\Theta(-z/E)/\prod_{j=1}^{l+1}(1+jz/E)$ with the power $v^{l+1}$. So the expansion is $\Theta(-z/E)\,v$ times a series of the shape \eqref{eq:set-S} with $r^{\mathrm{W}}_l$ in place of $r_l$. By \cref{prop:row-sp}\textup{(ii)} it solves the Picard--Fuchs equation, and the uniqueness statement after \cref{thm:row-pf}, with $r^{\mathrm{W}}_0=1$, gives $r^{\mathrm{W}}_l=r_l$. The identity is in the single rational variable $Y$, so it holds at all three roots. Hence $\redC(\Pn_0)=\hb\,\Srow$: the factor depending on $z$ only is determined, and the constants enter only through \cref{prop:gen-osc}.

\emph{The numerator and denominator step} \Lg{Lemma~5.9}. Write $\Tsum_g=\sum_k(\mathsf{p}_k/\mathsf{den})\,\sfF^k$, with $\mathsf{p}_k,\mathsf{den}\in\ringA$ and $\mathsf{den}\notin\primeq$; this is possible by \textup{(c)}. By \cref{thm:gen-main}, $\mathsf{den}\,\sF_g^{\Tmax}=\sum_k\mathsf{p}_kF^k$ in $\Kal\fps{y}$, and every term lies in the subring $\locOal\fps{y}$: the mirror map is independent of $\alpha$, $\GW_{g,d}^{\Tmax}\in\locOal$ by \cref{prop:gen-poly}, and $F\in\bbQ\fps{y}$. Applying the ring map $\redn$ gives $\redn(\mathsf{den})\,\sF_g^{\TF}=\sum_k\redn(\mathsf{p}_k)F^k$ in the domain $\bbK\fps{y}$, and $\redn(\mathsf{den})\neq0$ because $\mathsf{den}\notin\primeq$. Dividing in $\bbK\lser{y}$, and using the reduction of the trace and of the rows, gives \textup{(d)}. A limit would not be valid: at fixed $y_0$ the $y$-series of $Y_j(y)$ diverges for $\alpha$ near $\alpha^*$, because the branch points of $Y(y)$ tend to $y=0$ \Lg{Remark~5.10}.

\emph{Coefficient ring and transcendence.} The ring in \textup{(d)}, \textup{(e)} is obtained by an integrality argument. Finally, $F$ is transcendental over $\bbC(y)$ \Lg{Proposition~5.11}: $A$ has a logarithmic term at $\xi=0$, so its continuations $n$ times around $\xi=0$ are $A_n=A+2\pi\ri n\mathsf{b}\wcf$ with $\mathsf{b}\neq0$; if $F=\theta\log A$ were algebraic it would have finitely many branches, forcing $A_n=\mathsf{c}A_m$ for some $n\neq m$ and a constant $\mathsf{c}$, which is impossible.
